\documentclass[a4paper]{article}

\usepackage{graphicx}
\usepackage{amsmath,amssymb}
\usepackage{minted}
\usepackage{multirow}
\usepackage{float}
\usepackage{algorithm}
\usepackage{algpseudocode}
\usepackage{todonotes}
\usepackage{forest}
\usepackage{xstring}
\usepackage{xcolor}
\usepackage[most]{tcolorbox}
\usepackage{xparse}
\usepackage{natbib}

\title{iPad editable graphic companion persistence}
\date{2026-09-12}
\author{Mohammad Rahmani}

\begin{document}
\maketitle

\section{Problem}
Mind Mirror already stored PencilKit drawing data in a neighboring
\texttt{.ndgraphic} file, while LaTeX Source, Visual, and Preview used the raster
image. However, a newly associated ordinary raster initially received an empty
stored background. The bridge therefore had to infer the editable background
from the flattened raster when the document was reopened. In particular, a
legitimate white or transparent original background could later be mistaken for
an invalid legacy background after visible strokes had been flattened into the
paired image. Rename and move operations also moved both files but did not make
the sidecar's internal \texttt{image\_name} follow the raster's new basename.

\section{Fix}
\begin{itemize}
   \item The raster remains the canonical file referenced by LaTeX and displayed by Source, Visual, and Preview.
   \item The neighboring \texttt{.ndgraphic} remains the canonical editable iPad state containing PencilKit/vector data and canvas metadata.
   \item On the first iPad association, the companion now snapshots the exact original raster as its authoritative background before any flattened edit can replace the raster output.
   \item Version-2 companions carry an explicit \texttt{background\_is\_authoritative} marker, so valid white or transparent backgrounds are not replaced by a later flattened raster and editable strokes are not discarded.
   \item Untouched legacy companions can be migrated safely when they have no drawing/object data. Existing ambiguous legacy companions keep the previous recovery policy instead of being trusted blindly.
   \item Managed image/sidecar moves keep the two filenames paired and atomically rewrite the moved sidecar's internal \texttt{image\_name} to the new raster basename.
   \item Managed sidecar lifecycle recognition covers the raster formats already exposed by the Navigator iPad-edit command.
   \item Existing semantic source-reference/search path rules are unchanged: source references continue to target only the raster file, while the companion follows the raster internally.
\end{itemize}

\section{Regression coverage}
Regression checks verify first-association background persistence, preservation of
editable strokes over an authoritative white background, and metadata rebinding
when a raster/sidecar pair is renamed or moved.

\section{Repository hygiene}
The change extends the existing managed-graphic document model and filesystem
operation service. It does not add a second competing image workflow or a
runtime hotfix.

\end{document}
